\documentclass[12pt,a4paper,oneside]{article}
\usepackage[ngerman]{babel}
\usepackage[utf8]{inputenc}
\usepackage[babel,style=german]{csquotes}
\usepackage{amsmath}
\usepackage{amssymb}
\usepackage{amsthm}
\usepackage{enumitem}
\usepackage{xcolor}
\usepackage{empheq}
\usepackage{fancyhdr}
\usepackage{hyperref}
% \usepackage[draft]{hyperref}
\usepackage{icomma}
\usepackage{soul}
\setul{1.2pt}{0.5pt}
\usepackage{tikz}
\usepackage{tkz-euclide}
\usepackage{titlesec}
\usepackage{xurl}
\usepackage{geometry}
\geometry{
    left=2.28cm,
    right=2.28cm,
    top=3cm,
    bottom=3cm
}

\setlength{\parindent}{0pt}

\newcommand{\afz}[1]{\enquote{#1}}
\newcommand{\emptyline}{\hspace*{\fill}}
\newcommand{\N}{\mathbb{N}}
\newcommand{\R}{\mathbb{R}}

\begin{document}

    \section{Menge}

        \subsection{Grundlagen}

            Eine Menge ist ein mathematisches Modell, das bestimmte Objekte mit bestimmten Eigenschaften enthält.

            \vspace{-1em}
            \begin{align*}
                M = \{O \;\text{Objekt} \;|\; O \;\text{hat Eigenschaft}\; P\}
            \end{align*}

            Die Objekte können \textit{alles} sein: Zahlen, Formeln, Programme, Algorithmen, Dateien, Graphen, usw.
            Es ist ein ziemlich abstrakter Begriff.

            \emptyline

            Eine Menge kann auf diese Weise definiert werden:
            \afz{Es bezeichne $A$ die Menge, die die Zahlen 1, 2 und 3 enthält, und nichts sonst}.
            Und dann haben wir die Menge $A$:

            \vspace{-1em}
            \begin{align*}
                A = \{1, 2, 3\}
            \end{align*}

            Wir sagen: \afz{In $A$ ist die 42 nicht enthalten}.

            \emptyline

            Alle Objekte, die eine Menge enthält, heißen die Elemente der Menge.
            Seien $A$ eine Menge und $x$ ein Objekt.
            Wenn $x$ Element von $A$ ist, schreiben wir $x \in A$.
            Das Gegenteil wird als $x \notin A$ notiert.

            \emptyline

            Die Elemente einer Menge haben keine Reihenfolge,
            und jedes Element kommt in einer Menge höchstens einmal vor.
            Es gibt also keine Vielfachheiten, z.B.:

            \vspace{-1em}
            \begin{align*}
                \{1, 2, 3\} = \{3, 2, 1\} = \{1, 1, 2, 2, 3, 3, 3\}
            \end{align*}

            Die dritte Schreibweise nutzt man sehr selten, aber es kann vorkommen und ist mathematisch erlaubt.

            \emptyline

            Wir können eine Menge durch ihre enthaltenen Elemente eindeutig bestimmen.
            Seien $A$ und $B$ zwei Mengen.
            Dann gilt $A = B$, wenn für jedes Objekt $x$ gilt:

            \vspace{-1em}
            \begin{align*}
                x \in A \Longleftrightarrow x \in B
            \end{align*}

        \subsection{Teilmengen}

            Eine Menge kann Teilmenge einer anderen Menge sein.
            Andererseits kann eine Menge auch Obermenge einer anderen Menge sein.

            \emptyline

            Seien $A$ und $B$ zwei Mengen.
            Wenn $A$ Teilmenge von $B$ ist, schreiben wir $A \subseteq B$.
            In diesem Fall ist $B$ Obermenge von $A$, geschrieben wird $B \supseteq A$.
            Folglich wenn $A \subseteq B$ und $B \subseteq A$ gleichzeitig gelten, dann gilt $A = B$.

            \emptyline

            Wir merken, dass das Teilmengensymbol $\subseteq$ eine Linie am unteren Teil hat.
            Das bedeutet, dass es auch $A = B$ sein kann, wenn $A$ Teilmenge von $B$ ist.
            Dazu ergänzen wir den Begriff \textit{echte Teilmenge}.
            Die sind die Teilmengen, die nicht identisch mit ihrer Obermenge sind.
            Dafür gibt es verschiedene Schreibweisen: $A \subsetneqq B$, $A \subsetneq B$ oder einfach $A \subset B$.

        \subsection{Leere Menge}

            Die leere Menge ist die Menge, die gar keine Elemente enthält.
            Dafür gibt es verschiedene Schreibweisen: $\{\}$, $\emptyset$ und $\varnothing$.

            \emptyline

            Die leere Menge ist Teilmenge jeder Menge.
            Aber man muss beachten, dass $\emptyset$ und $\{\emptyset\}$ ungleich sind.
            Ersteres beschreibt die leere Menge selbst, Letzteres beschreibt eine Menge,
            die die leere Menge als Element enthält.

        \subsection{Definitionsweise einer Menge}

            Es gibt verschiedene mögliche Weisen, eine Menge zu definieren.
            Aber wir müssen beachten, dass eine Definition exakt und nicht mehrdeutig sein muss.

            \emptyline

            Es gibt ein Beispiel, für ein $n \in \N_0$ definiere:

            \vspace{-1em}
            \begin{align*}
                [n] = \{1, \cdots, n\}
            \end{align*}

            Aber wir können (an der Uni) nicht direkt behaupten, was die Punkte bedeuten.
            Obwohl es so aussieht, dass die Punkte alle natürlichen Zahlen zwischen 1 und $n$ repräsentiert,
            ist diese Definition trotzdem mehrdeutig.
            
            \emptyline

            Deswegen definieren wir $[n]$ lieber wie:

            \vspace{-1em}
            \begin{align*}
                [n] = \{i \in \N_0 \;|\; 1 \leq i \leq n\}
            \end{align*}

            und ergänzen wir noch einen speziellen Fall, damit für alle mögliche $n$ eine entsprechende Menge $[n]$ gibt:

            \vspace{-1em}
            \begin{align*}
                [0] = \emptyset
            \end{align*}

            Dazu kommen wir zu \textit{set comprehension}.
            Sei $P(x)$ eine Aussage für jedes Objekt $x$ wahr oder falsch ist.
            Dann enthält $\{x \in M \;|\; P(x)\}$ genau die $x \in M$, für die $P(x)$ wahr ist.

            \emptyline

            Es ist nicht notwendig, in die Definition einen Definitionsbereich für $x$ einzusetzen.
            Falls das auch verständlich ist, können wir natürlich auch $\{x \;|\; P(x)\}$ schreiben.
            Es gibt auch eine alternative Schreibweise $\{x : P(x)\}$.
            Der Zweck ist nur, wie gesagt, die Definition exact zu sein.

            \emptyline

            \textbf{ABER} es ist notwendig, nur \afz{harmlose} Aussagen $P$ zu nutzen (vgl. Russells Paradoxon).
            Eine Menge $A = \{x \;|\; x \notin A\}$ oder $R = \{x \;|\; x \notin x\}$ führt zu einem Widerspruch.

        \subsection{Mengenoperatoren}

            Seien $A$ und $B$ zwei Mengen.

            \emptyline

            Dann ist die \ul{Vereinigung} von $A$ und $B$:

            \vspace{-1em}
            \begin{align*}
                A \cup B = \{x \;|\; x \in A \vee x \in B\}
            \end{align*}

            \newpage

            Wenn ein Objekt $x$ ein Element von $A$ und $B$ ist, zählt das in $A \cup B$ auch nur einmal.
            Weil wie gesagt, ein Element kommt in einer Menge höchstens einmal vor, z.B.:

            \vspace{-1em}
            \begin{align*}
                \{1, 2, 3\} \cup \{3, 4, 5\} = \{1, 2, 3, 4, 5\}
            \end{align*}

            Der \ul{Durchschnitt} von $A$ und $B$ ist so definiert:

            \vspace{-1em}
            \begin{align*}
                A \cap B = \{x \;|\; x \in A \wedge x \in B\}
            \end{align*}

            Dazu führen wir einen neuen Begriff ein:
            \ul{disjunkte Mengen} sind die Mengen $A$ und $B$, für die $A \cap B = \emptyset$ gilt.
            Wir sagen in diesem Fall, dass $A$ und $B$ disjunkt sind.

            \emptyline

            Die \ul{Mengendifferenz} von $A$ und $B$, gelesen \afz{$A$ ohne $B$}, ist so definiert:

            \vspace{-1em}
            \begin{align*}
                A \setminus B = \{x \in A \;|\; x \notin B\}
            \end{align*}

            Die \ul{Kardinalität} einer Menge $A$ ist die Anzahl ihrer Elemente (hier nur für endliche Mengen definiert),
            notiert als $|A|$ (evtl. $\#A$ oder $\text{card}A$).
            Die Kardinalität für eine endliche Menge ist eine nichtnegative ganze Zahl.
            Bei unendlichen Mengen reden wir hier nicht über ihre Kardinalität (aber man kann das ordentlich definieren).

            \emptyline

            Man muss achten, dass die Summe von zwei Kardinalitäten zweier Mengen nicht unbedingt gleich wie
            die Kardinalität der Vereinigung zweier Mengen, z.B.:

            \vspace{-1em}
            \begin{align*}
                |\{1, 2\}| + |\{2, 3\}| = 2 + 2 = 4 \neq 3 = |\{1, 2, 3\}| = |\{1, 2\} \cup \{2, 3\}|
            \end{align*}

            Letztendlich ist die \ul{Potenzmenge} einer Menge $M$:

            \vspace{-1em}
            \begin{align*}
                \mathcal{P}(M) = 2^M = \{A \;|\; A \subseteq M\}
            \end{align*}

            Hier ist $2^M$ eine spezielle Schreibweise. 

    \newpage

    \section{Abbildung}

        Seien $A$ und $B$ zwei (nicht notwendigerweise verschiedene) Mengen.

        \emptyline

        Dann ordnet eine Abbildung $f: A \to B$ jedem Element von $A$ genau ein Element von $B$ zu.
        Hier ist $A$ der Definitionsbereich von $f$, $B$ der Zielbereich von $f$.

        \emptyline

        Eine Abbildung ist bestimmt durch Abbildungsvorschrift, z.B. Spezifikation von Funktionswerten.
        Wird ein Element $x \in A$ auf das die Abbildung angewendet wird,
        dann wird $x \in A$ auf $f(x) \in B$ abgebildet, notiert als:

        \vspace{-1em}
        \begin{align*}
            x \mapsto f(x)
        \end{align*}

        Abbildungen heißen auch Funktionen, aber Funktionen beschreiben meist dem Fall, wenn Definitions- und Zielbereich Zahlenmengen sind.
        Abbildungen sind normaleweise abstrakter.

        \emptyline

        Wir stellen ein Beispiel vor (Collatz Problem):

        \vspace{-1em}
        \begin{align*}
            f: \N_0 \to \N_0, n \mapsto
            \begin{cases}
                n / 2, &\text{falls $n$ gerade} \\
                3n + 1, &\text{falls $n$ ungerade}
            \end{cases}
        \end{align*}

        Oder:

        \vspace{-1em}
        \begin{gather*}
            \psi: \{G \;|\; G \;\text{Graph}\} \to \{0, 1\}, \pi: \{G \;|\; G \;\text{Graph}\} \to \N_0 \\
            \begin{aligned}
                \psi(G) &=
                \begin{cases}
                    1, &\text{falls $G$ eine Euler-Tour hat} \\
                    0, &\text{sonst}
                \end{cases} \\
                    \pi(G) &= \text{Anzahl der Euler-Touren von $G$}
            \end{aligned}
        \end{gather*}

        Aus diesem Beispiel können wir erkennen, dass die Eingaben einer Abbildung nicht unbedingt Zahlen sind.

        \emptyline

        Zu jeder Abbildung $f$ gehört Definitionsbereich, Zielbereich und Abbildungsvorschrift.
        Wenn der Definitionsbereich oder Wertebereich zweier Abbildungen unterschiedlich sind, dann sind sie verschiedene Abbildungen.
        Aber wenn nur die Abbildungsvorschrift unterschiedlich aussieht, kann das sein, dass die Abbildungen trotzdem gleich sind, z.B.:

        \vspace{-1em}
        \begin{align*}
            &f_1: \R \to \R, x \mapsto x^2 \\
            &f_2: \R \to \R_+, x \mapsto x^2 \\
            &f_3: \R \to \R, x \mapsto x(x + 1) - x
        \end{align*}

        Hier sind Abbildungen $f_1$ und $f_3$ gleich, aber $f_1$ und $f_2$ sind verschiedene Abbildungen.

\end{document}
